\subsection{Phänomenologische Brücke: Kepler-Anker, 120-Zeller und Re/Ra-Dualität}\label{sec:eabc-phaenomenologie}

Die in den Abschnitten~\ref{sec:determ-numerik-fixpunkte}--\ref{sec:eabc-zeugen-phi} dokumentierten numerischen Läufe schließen eine Lücke zwischen arithmetischer EABC-Integrität und klassischen Grenzwerten der Kontinuumsphysik. Am Kepler-Anker $M=113160$---dem kombinatorischen Fixpunkt der H32-Phasenschale und des quaternionischen Pfad-B-Filters (Abschnitt~\ref{sec:empirie-resonanz})---zerfallen die Hurwitz-Einheiten des regulären 24-Zells in zwei disjunkte Zeugenklassen: rationale Bauten auf der reellen $e$-Achse ($a=0$) und algebraische $\Phi$-Kopplungen ($a\neq 0$), deren Urzeuge $(0,-1,0,0)$ die ikosaedrische Kernsymmetrie trägt.

\paragraph{Rationaler Sektor und Impulsdissipation.}
Die rationale Kaskade $M/\mathrm{Norm}\in\{113160,56580,37720,\ldots,18860\}$ ordnet ganzzahlige Schalenverhältnisse der Gitternorm einer effektiven Reynolds-Zahl zu (Tabelle~\ref{tab:eabc-rational}). Die Eichung $\Lambda_{\mathrm{impuls}}$ an der maximal verdichteten Schale $\mathrm{Norm}=6$ fixiert den laminar-turbulenten Umschlag bei $Re=2300$; höhere harmonische Stufen ($13800$, $6900$, $4600$) interpretieren wir \emph{hypothetisch} als stabil-laminare Resonanzüberbauten oberhalb der viskosen Grenze. Diese Zuordnung ist keine Ableitung der Navier--Stokes-Gleichungen aus der Hurwitz-Arithmetik, sondern eine \emph{dimensionslose Kalibrierung} zwischen diskreter $M$-Kaskade und einem etablierten Fluiddynamik-Schwellenwert.

\paragraph{$\Phi$-Sektor und thermische Konvektion.}
Dual dazu koppelt der $\Phi$-Sektor die algebraischen Kerne $\alpha+\beta\Phi\in\mathbb{Z}[\Phi]$ an effektive Rayleigh-Zahlen (Tabelle~\ref{tab:eabc-phi}). Der Urzeuge mit $Ra=1708$ markiert den Einstieg thermischer Konvektionswalzen (Bénard-Instabilität); die absteigenden Stufen fallen in exakten rationalen Verhältnissen $\tfrac{3}{10}$, $\tfrac{3}{35}$, $\tfrac{1}{4}$, $\tfrac{3}{40}$, $\tfrac{1}{24}$ relativ zu $Ra_{\mathrm{krit}}$ ab. Im Forschungsprogramm lesen wir dies als Hinweis, dass die $\mathbb{Z}[\Phi]$-Integritätsbedingung die Energiedissipation in quantisierten Schritten entlang der fünffachen Symmetrie des Ikosaeders bzw.\ der dichten Packung des 120-Zells (H4-Komplex in $\mathbb{R}^4$, verwandt mit Keplers Kugelgeometrie) strukturiert.

\paragraph{Synthese.}
Die Dualität Re/Ra spiegelt die in Axiom~\ref{ax:split} vorgeschlagene Zerlegung der $E_8$-Wurzeln in Torsions- ($\mathcal{S}_{\mathbb{Z}}$, $112$ Vektoren) und Translationsschale ($\mathcal{S}_{\mathbb{Z}+\frac12}$, $128$ Vektoren): Impulsdominierte rationale Fixpunkte auf der einen, thermisch-algebraische $\Phi$-Zeugen auf der anderen Seite. Eine kausale Verknüpfung mit Orbitaloperatoren $\Phi_T$, $\Phi_R$ (Axiom~\ref{ax:orbit}) oder mit Vermutung~\ref{conj:RH} bleibt offen; die hier berichteten Tabellen liefern jedoch reproduzierbare, widerlegbare Kalibrierungsdaten für künftige Spezifikation dieser Abbildungen.

\begin{remark}[Methodische Einordnung]
Die Re/Ra-Zuordnung ist eine \emph{phänomenologische Heuristik} innerhalb des \#Energiedoku-Programms. Klassische Schwellenwerte ($Re\approx 2300$, $Ra\approx 1708$) dienen als externe Kalibrierungsanker; exakte Übereinstimmung an den gewählten Symmetriepunkten folgt aus der Eichungsdefinition, nicht aus einem unabhängigen physikalischen Beweis.
\end{remark}
